% tex/math.tex — Math symbols, operators, and shared notation
%
% NOTE: theorem/lemma/definition environments are defined in `tex/env.tex` to
% keep all environment configuration in a single place, aligned with the
% reference material under `tex/literature/`.

%==============================================================================
% MATH OPERATORS
%==============================================================================
\providecommand{\argmax}{\mathop{\rm arg\,max}}
\providecommand{\argmin}{\mathop{\rm arg\,min}}
\providecommand{\sign}{\mathop{\rm sign}}
\providecommand{\Tr}{\mathop{\rm Tr}}

%==============================================================================
% PROBABILITY AND STATISTICS
%==============================================================================
\providecommand{\E}{\mathbb{E}}
\providecommand{\Var}{\mathrm{Var}}
\providecommand{\Cov}{\mathrm{Cov}}
\providecommand{\R}{\mathbb{R}}

%==============================================================================
% POLICY NOTATION
%==============================================================================
\providecommand{\piroll}{\pi_{\mathrm{roll}}}
\providecommand{\pitheta}{\pi_{\theta}}
% PPO-style implementations often denote the rollout snapshot as \piold. In this
% paper, the theoretical behavior/data-generating policy is always \piroll; we
% only mention \piold when mapping to implementation terminology.
\providecommand{\piold}{\pi_{\mathrm{old}}}
\providecommand{\piref}{\pi_{\mathrm{ref}}}
\providecommand{\ylt}{y_{<t}}
\providecommand{\yle}{y_{\le t}}

%==============================================================================
% IMPORTANCE RATIOS (TRM/GSPO STYLE)
%------------------------------------------------------------------------------
% We follow the reference style in `tex/literature/` and use:
% - token-level ratio: \rho_t
% - sequence-level geometric-mean ratio: \rhoseq
% - prefix-level geometric-mean ratio: \rhoprefix{t}
%==============================================================================
\providecommand{\rhoseq}{\rho^{\mathrm{seq}}}
\providecommand{\rhoprefix}[1]{\rho_{#1}^{\mathrm{pre}}}

%==============================================================================
% DIVERGENCE MEASURES (following TRM paper convention)
%==============================================================================
\providecommand{\DKL}{D_{\mathrm{KL}}}
\providecommand{\DTV}{D_{\mathrm{TV}}}
\providecommand{\Dkltok}{D_{\mathrm{KL}}^{\mathrm{tok}}}
\providecommand{\Dtvtok}{D_{\mathrm{TV}}^{\mathrm{tok}}}
\providecommand{\Dkltokmax}{D_{\mathrm{KL}}^{\mathrm{tok,max}}}
\providecommand{\Dtvtokmax}{D_{\mathrm{TV}}^{\mathrm{tok,max}}}
\providecommand{\Dklseq}{D_{\mathrm{KL}}^{\mathrm{seq}}}
\providecommand{\Dtvseq}{D_{\mathrm{TV}}^{\mathrm{seq}}}
\providecommand{\Dbar}{\bar{D}}

%==============================================================================
% COMMON MATH SHORTCUTS
%==============================================================================
\providecommand{\Ls}{\mathcal{L}}
\def\eps{{\epsilon}}
\providecommand{\epslow}{\epsilon_{\mathrm{low}}}
\providecommand{\epshigh}{\epsilon_{\mathrm{high}}}
\def\ceil#1{\lceil #1 \rceil}
\def\floor#1{\lfloor #1 \rfloor}
